\section{Introduction}

Causal discovery — the inference of cause-and-effect relationships from purely observational data — is a central challenge in applied machine learning across disciplines as diverse as neuroscience, biology, economics, and finance. In these domains, understanding the causal structure of a system is crucial for making reliable predictions under interventions, designing robust policies, and explaining observed phenomena beyond mere correlation. Formally, the objective is to recover a Directed Acyclic Graph (DAG) $\mathcal{G} = (\mathcal{V}, \mathcal{E})$ over the set of observed variables $\mathcal{V}$, where a directed edge $(i \rightarrow j) \in \mathcal{E}$ encodes the hypothesis that variable $i$ is a direct Granger cause of variable $j$.

\subsection{The Nonstationarity Challenge}

In temporally ordered multivariate datasets, the challenge is compounded by the presence of lag-dependent interactions. Temporal causal discovery methods operate on the principle of Granger causality \cite{granger1969investigating}: variable $X^i$ is said to Granger-cause $X^j$ if the past values of $X^i$ improve predictions of the future of $X^j$ beyond what is achievable from the history of $X^j$ alone. Classical algorithms such as the PC algorithm and its temporal extension PCMCI+ \cite{runge2019detecting} rely on conditional independence testing to identify Granger-causal relationships, while score-based approaches such as DYNOTEARS \cite{pamfil2020dynotears} frame the problem as a continuous optimization.

However, all of these methods share a critical and often unexamined assumption: the data-generating process is \textit{stationary}. In financial markets, this assumption is routinely violated. Stock price correlations, cross-asset spillover effects, and inter-market transmission mechanisms shift dramatically across different market regimes — from bull markets to bear markets, from periods of low volatility to crisis-driven turbulence. These structural breaks render any single-model causal estimate invalid across the entire time span. An algorithm trained on pre-crisis dynamics will output a causal graph that is not only incorrect but potentially misleading in its post-crisis predictions.

\subsection{Neural Approaches and their Limitations}

Recent advances in deep learning have produced powerful neural architectures for causal discovery. CUTS+ \cite{cheng2023cuts} deploys a coarse-to-fine Message-Passing GNN to handle high-dimensional, irregular time series. NoCurl and NOTEARS \cite{Zheng2018} introduced the foundational continuous differentiable DAG constraint $h(\mathbf{W}) = \text{tr}(e^{\mathbf{W} \circ \mathbf{W}}) - N = 0$, enabling gradient-based optimization over the combinatorially intractable DAG space. These methods demonstrate strong performance on stationary benchmarks but are fundamentally agnostic to temporal nonstationarity. When applied to multi-regime data, they produce globally smoothed, blurred causal estimates that fail to capture the true regime-dependent structure.

\begin{revenv}
While some existing methods attempt to model nonstationarity directly, they exhibit critical shortcomings in complex financial scenarios. Algorithms such as CD-NOD \cite{huang2020causal} and Regime-PCMCI \cite{saggioro2020causal} introduce regime-awareness, but they often struggle with the high dimensionality and irregular sampling frequencies typical of real-world financial data. More importantly, these purely data-driven nonstationary methods lack a mechanism to incorporate domain-specific topological priors.

Furthermore, neural causal discovery is notoriously sample-inefficient. The combinatorial complexity of the DAG search space grows super-exponentially with the number of variables. In sparse financial datasets without strong domain knowledge, algorithms frequently converge to suboptimal local minima, particularly when the time series is short within a specific turbulent market regime.
\end{revenv}

\subsection{The REIGN Framework}

To address these fundamental limitations, we propose the \textbf{Regime-Enhanced Intelligent Granger Network (REIGN)}, an end-to-end causal discovery framework engineered from the ground up to handle nonstationary environments. REIGN's architecture integrates four synergistic components. First, a robust preprocessing pipeline applies MICE-based imputation and outlier handling. Second, the PELT algorithm detects changepoints and segments the time series into locally stationary regimes, ensuring each downstream model operates on a homogeneous distribution. Third, a novel prior-injection mechanism queries a Large Language Model (LLM) to extract domain-specific zero-shot topological priors, dramatically constraining the neural search space without requiring labeled training data. Finally, a per-regime MPGNN engine optimizes a continuous DAG-constrained adjacency matrix, anchored by the LLM prior, before a confidence-weighted ensemble aggregates the local models into a globally consistent causal graph.

The primary contributions of this work are four-fold:
\begin{enumerate}
    \item We introduce a principled architecture for regime-aware temporal causal discovery, employing the PELT algorithm to partition nonstationary time series into stationary local windows and enabling each neural Granger model to specialize on a homogeneous distributional regime.
    \item We pioneer the integration of zero-shot LLM-generated structural priors into a differentiable continuous Lagrangian optimization objective, bridging the gap between human domain expertise and deep causal learning without requiring explicit human annotation.
    \item We propose a novel confidence-weighted ensemble mechanism that synthesizes regime-specific local DAGs into a globally consistent summary graph through temporal importance weighting.
    \item We demonstrate through rigorous empirical evaluation on synthetic nonstationary VAR datasets that REIGN achieves state-of-the-art performance, outperforming established baselines such as PCMCI+ by over 11\% in F1 score and dominating across all continuous structural metrics.
\end{enumerate}
